\documentclass[11pt,letterpaper]{article}
\usepackage[margin=0.60in,top=0.58in,bottom=0.56in]{geometry}
\usepackage{mathptmx}
\usepackage{microtype}
\usepackage{graphicx}
\usepackage{booktabs}
\usepackage{array}
\usepackage{longtable}
\usepackage{amsmath,amssymb}
\usepackage[hidelinks]{hyperref}
\setlength{\parindent}{0pt}
\setlength{\parskip}{2.5pt}
\setlength{\columnsep}{0.25in}
\setcounter{secnumdepth}{0}
\renewcommand{\section}[1]{\vspace{3pt}\noindent\textbf{\underline{#1}}\par\vspace{1pt}}
\renewcommand{\subsection}[1]{\vspace{2pt}\noindent\textbf{#1}\par\vspace{1pt}}
\renewcommand{\familydefault}{\rmdefault}
\providecommand{\tightlist}{\setlength{\itemsep}{0pt}\setlength{\parskip}{0pt}}
\providecommand{\pandocbounded}[1]{\sbox0{#1}\ifdim\wd0>\linewidth\resizebox{\linewidth}{!}{#1}\else\usebox0\fi}
\title{\vspace{-1.2em}\textbf{\MakeUppercase{Neurologically Motivated
Self-Supervised Learning for Bilateral Gait Geometry}}\vspace{-0.5em}}
\author{\normalsize [Authors and affiliations to be inserted in the RISEx template]}
\date{}
\begin{document}
\twocolumn[
  \begin{@twocolumnfalse}
  \maketitle
  \vspace{-1.1em}
  \end{@twocolumnfalse}
]
\subsection{INTRODUCTION}\label{introduction}

Human-aware robots often begin with a perception module that summarizes
a person's pose and motion. Before such a summary is used for prediction
or interaction, it should be tested for geometric information relevant
to the task. We study a small but concrete case: a reflected skeleton
swaps anatomical sides and reverses a left--right speed contrast. Stroke
and Parkinson's disease can alter gait asymmetry {[}1{]}, which
motivates the geometry; our target remains a pose-derived movement score
rather than a clinical label. We test whether masked joint-embedding
predictive architecture (JEPA) training makes this score easier to
recover.

\subsection{MATERIALS AND METHODS}\label{materials-and-methods}

GAVD contributes 625 quality-checked clips from 93 source videos
{[}2{]}. We process each clip independently, retaining its 33 landmarks
in a 64-time-step, three-coordinate array. Four consecutive steps form
one patch, preserving a 16-by-33 time-by-joint layout. The online
encoder sees visible patches, a predictor estimates hidden-patch
features, and an exponential-moving-average teacher provides targets
{[}3{]}. No movement score or affected-side label enters encoder
training.

The target is the mean normalized difference of left and right median
speeds across five paired joints. Five outer folds hold out complete
source videos; every clip from a held-out source remains outside encoder
training, readout fitting, and readout tuning. Five seeds repeat the
matched comparison of each trained encoder with its own starting
weights. We resample whole sources for paired intervals while holding
fitted models fixed.

\subsection{RESULTS AND DISCUSSION}\label{results-and-discussion}

The initial encoder reached \(R^2=0.223\) with a motion-sensitive frozen
readout. The five trained teacher encoders reached \(0.101\)--\(0.114\);
their paired differences from initialization ranged from \(-0.109\) to
\(-0.122\). Yet trained predictors preferred correct-clip teacher
features to features from another source in 375 of 375 diagnostic
checks, compared with 33 of 75 checks before training. Feature
correspondence therefore improved without improving the tested
bilateral-motion estimate.

This is a useful failure mode for human-aware perception, but its reach
is limited. The expanded readout also adds temporal and missingness
statistics, and preprocessing changes the target calculation. The
evidence identifies a shortfall in this training-and-readout pipeline
rather than proving that trained features lack all geometric motion
information.

\subsection{CONCLUSIONS}\label{conclusions}

Masked feature prediction can capture clip-related structure while
making a relevant bilateral motion quantity less accessible to a frozen
linear readout. The result motivates explicit task-level checks before
using self-supervised articulated-motion features downstream. It does
not claim robotic performance or control benefit.

\subsection{REFERENCES}\label{references}

{[}1{]} Lewek MD et al.~\emph{Gait Posture} 31:256--260, 2010. {[}2{]}
Ranjan R et al.~\emph{IEEE Access} 13:45321--45339, 2025. {[}3{]} Assran
M et al.~arXiv:2301.08243, 2023.
\end{document}
